\section{Preparing infinite factors and comparing modular graphs}
\label{sec:modular-preparation}

We now pass from infinite factors to a common modular setting.  The input is
ordinary closeness on a separable Hilbert space.  After one small adjustment,
proper infiniteness makes that closeness completely bounded.  A fixed
hyperfinite tensor factor then supplies an irreducible finite subfactor in a
state centralizer.  Aligning its left and right actions gives a common cyclic
separating vector.  The main estimate of this section compares the two Tomita
graphs uniformly under rescaling.  In \Cref{sec:modular-alignment}, this
estimate will permit a small change of position after which the modular
group of one algebra preserves both algebras.

\subsection{A fixed tensor factor and common position}

\begin{lemma}[Amplification by common isometries]
\label{lem:common-infinite-amplification}
Let $M,N\subseteq\B(H)$ contain the same unital copy of $\B(\ell^2)$.
Then $\dcb(M,N)=d(M,N)$.
\end{lemma}

\begin{proof}
For each $n\geq1$ choose isometries $s_1,\ldots,s_n$ in that common copy
such that $s_i^*s_j=\delta_{ij}1$ and $\sum_i s_is_i^*=1$.
The operator $V:H^n\to H$ given by $V(\xi_i)=\sum_i s_i\xi_i$ is unitary.
For either $P=M$ or $P=N$,
\[
 V M_n(P)V^*=P,
 \qquad V[x_{ij}]V^*=\sum_{i,j}s_ix_{ij}s_j^*.
\]
Indeed the inverse coefficient formula is $x_{ij}=s_i^*xs_j$.
The same unitary therefore identifies the two amplified unit balls with the
two original unit balls.  Taking the supremum over $n$ proves the assertion.
\end{proof}

We fix, once and for all,
\[
 \lambda=e^{-1},\qquad \mu=e^{-\sqrt2}.
\]
For $0<\nu<1$, let $\rho_\nu$ be the state of $M_2(\C)$ with density
$(1+\nu)^{-1}\operatorname{diag}(1,\nu)$.  Form a countable infinite
tensor product containing infinitely many sites of each of the two types
$(M_2(\C),\rho_\lambda)$ and $(M_2(\C),\rho_\mu)$, and denote its von
Neumann algebra and product state by $(D,\rho)$.  We use the tensor product
of the site standard representations, on a separable Hilbert space $K$.
The product vector $\Omega_\rho$ is cyclic for both left and right local matrix actions,
and hence is cyclic and separating for $D$.  Thus $\rho$ is faithful and
normal.  Both $D$ and its standard commutant are hyperfinite: their finite
site matrix algebras increase strongly to them.

We recall the modular notation. In a standard representation with a
cyclic separating vector $\Omega$ implementing a faithful normal state
$\psi$, the closure of $x\Omega\mapsto x^*\Omega$ has polar
decomposition $S=J\Delta^{1/2}$. The antilinear isometry $J$ is the
modular conjugation, and $\sigma_t^\psi(x)=\Delta^{it}x\Delta^{-it}$
is the modular automorphism group. The state determines this action
independently of the chosen standard representation; see
\cite{Takesaki2003II} for Tomita--Takesaki theory.

For a faithful normal state $\psi$ on an algebra $P$, its
\emph{centralizer} is
\[
 P_\psi=\{a\in P:\psi(ax)=\psi(xa)\text{ for every }x\in P\}.
\]
It is also the fixed algebra of the modular automorphism group
$\sigma^\psi$.  We use the following precise continuous-core facts.  The
modular crossed product $P\rtimes_{\sigma^\psi}\R$ is semifinite,
with a faithful normal semifinite trace scaled by its
dual action \cite[III.4.8.1--4, p.~329]{BlackadarOperatorAlgebras}.
For a type $\III$ factor $P$, this crossed product is a factor exactly
when $P$ is of type $\III_1$; in that case it is of type $\II_\infty$
\cite[Lemma~8.2 and Corollary~9.7]{Takesaki1973}.
For a semifinite factor, choose a faithful normal semifinite trace $\tau$.
The modular crossed product is independent, up to isomorphism, of the
faithful normal semifinite weight used to form it
\cite[III.4.8.1, p.~329]{BlackadarOperatorAlgebras}.
Since $\sigma^\tau$ is trivial, this crossed product is
$P\barotimes L(\R)\cong P\barotimes L^\infty(\R)$ and is not a factor.
Here $L(\R)$ is the von Neumann algebra generated by translations on
$L^2(\R)$, identified with $L^\infty(\R)$ by Fourier transform.
Thus a factor has a factorial modular crossed product precisely when
it is of type $\III_1$.

\begin{lemma}[The fixed tensor factor]
\label{lem:fixed-tensor-factor}
One has $D_\rho'\cap D=\C1$.  If $A$ is any factor with separable predual,
then $A\barotimes D$ is a type $\III_1$ factor and
\[
 A\barotimes D\cong(A\barotimes D)\barotimes R_\lambda,
\]
where $R_\lambda$ is the Powers factor obtained using only
$(M_2(\C),\rho_\lambda)$ sites.
\end{lemma}

\begin{proof}
Swapping two sites of the same type is implemented by a unitary in a finite
site algebra.  This unitary preserves the product state, so it belongs to
$D_\rho$.  Suppose $x\in D_\rho'\cap D$.  Given a finite site element $a$,
choose bounded finite site approximants $x_j$ with
$\norm{(x-x_j)\Omega_\rho}\to0$.  A product of the indicated swaps moves
$a$ to an element $a_j$ supported away from $x_j$.  Since $x$ commutes with
the implementing unitary and the state is invariant,
\[
 \rho(ax)=\rho(a_jx),\qquad
 \rho(a_jx_j)=\rho(a)\rho(x_j).
\]
State Cauchy--Schwarz bounds the error in the first product by
$\rho(aa^*)^{1/2}\norm{(x-x_j)\Omega_\rho}$; the other error tends to zero
as well.  Hence $\rho(ax)=\rho(a)\rho(x)$ for every finite site $a$.
Density gives $(x-\rho(x)1)\Omega_\rho=0$, and faithfulness gives
$x=\rho(x)1$.  In particular $D$ is a factor.

Choose a faithful normal state $\psi$ on $A$ and put
$P=A\barotimes D$ with state $\psi\otimes\rho$.  In the regular
representation of its modular crossed product on $L^2(\R,H_P)$, the
generators are
\[
 (\pi(x)\xi)(t)=\sigma^{\psi\otimes\rho}_{-t}(x)\xi(t),
 \qquad (\lambda_s\xi)(t)=\xi(t-s).
\]
Here $H_P$ can be the standard space of $P$.  Let $Z$ be central in the
crossed product.  After Fourier transform in $t$, commutation with every
$\lambda_s$ makes $Z$ a decomposable operator.  The crossed product also
commutes with the constant representation of $P'$.  Consequently its
decomposable field takes values in $P$.  Commutation with the constant
operators $1\otimes D_\rho$ puts this field in $A\otimes1$: normal slices
and $D_\rho'\cap D=\C1$ prove this relative-commutant identity.

An off-diagonal matrix unit at a $\nu$-site is a modular eigenoperator
with frequency $\log\nu$ or $-\log\nu$.  Its regular representation,
after Fourier transform, is that matrix unit followed by a frequency
translation.  If $a(\zeta)\otimes1$ denotes the field of $Z$, commutation
therefore gives
\[
 a(\zeta+\log\lambda)=a(\zeta),\qquad
 a(\zeta+\log\mu)=a(\zeta)
 \quad\text{almost everywhere}.
\]
The generated subgroup is countable and dense in $\R$.  Test the field on
countably many vector coefficients.  Translation continuity in local
$L^1$ extends its invariance to every real translation, and convolution
with compactly supported approximate identities shows that each
coefficient is constant.  Thus $Z=a_0\otimes1$ is a constant field.
Commutation with $\pi(A\otimes1)$, using strong continuity of the modular
action and a countable strongly dense set, gives $a_0\in A'\cap A=\C1$.
The core is a factor, so the continuous-core criterion makes $P$ a type
$\III_1$ factor.

Finally, adjoining another countable family of $\lambda$-sites and
reindexing that family gives the asserted absorption.  All these tensor
products preserve the displayed product states.  The usual identification
of the two-ratio product with the hyperfinite $\III_1$ factor is also
recorded in \cite[Example~1.6, p.~103]{Haagerup1987}; the argument above
establishes the additional assertion for every $A$.
\end{proof}

For a faithful normal state $\psi$ on $P$, its \emph{bicentralizer}
$\mathrm{B}(P,\psi)$ consists of those $a\in P$ for which
$\norm{[a,x_n]}_\psi\to0$ whenever $(x_n)$ is a bounded sequence in $P$
and $\norm{x_n\psi-\psi x_n}\to0$.  Here
$\norm{x}_\psi=\psi(x^*x)^{1/2}$, and the last norm is the predual norm
of the functional $y\mapsto\psi(yx_n-x_ny)$.
We use the following consequence of two established results.  If a
$\sigma$-finite type $\III_1$ factor $P$ absorbs $R_\lambda$ for some $0<\lambda<1$, then
its bicentralizer is $\C1$ \cite[Theorem~D(ii)]{Marrakchi2020}.  For a
type $\III_1$ factor with separable predual, trivial bicentralizer gives an
irreducible hyperfinite $\II_1$ subfactor $R\subseteq P$ with a faithful
normal conditional expectation $E:P\to R$.  The finite subfactor is
constructed in \cite[Lemma~3.8 and the proof of Theorem~3.1,
$(2)\Rightarrow(3)$, pp.~142--145]{Haagerup1987}: the construction first gives such a subfactor
in a nonzero corner, which is isomorphic to $P$.  This statement imposes
no injectivity assumption on $P$.  The state $\tau_R\circ E$ centralizes
$R$, by bimodularity of $E$.

\begin{definition}[Prepared pair]
\label{def:prepared-pair}
A \emph{prepared pair} consists of type $\III_1$ factors
$M,N\subseteq\B(H)$ on a separable Hilbert space, a common cyclic
separating unit vector $\Omega$, and a hyperfinite $\II_1$ factor
$R\subseteq M\cap N$ such that
\begin{enumerate}[label=\textup{(\roman*)}]
\item $R'\cap M=R'\cap N=\C1$;
\item the vector state $\varphi(x)=\ip{x\Omega}{\Omega}$ centralizes $R$
when restricted to either algebra;
\item $J_MrJ_M=J_NrJ_N$ for every $r\in R$, where $J_M,J_N$ are the
modular conjugations associated with $\Omega$.
\end{enumerate}
\end{definition}

\begin{proposition}[Preparing an infinite pair]
\label{prop:prepared-pair}
There are absolute constants $c,C>0$ with the following property.
Suppose $M_0,N_0\subseteq\B(H)$ are factors on a separable
Hilbert space, $M_0$ is infinite, and $d(M_0,N_0)<t<c$.
For the fixed $D\subseteq\B(K)$
above, there are unitaries $u_0\in\B(H)$ and
$v\in\B(H\otimes K)$ satisfying
\[
 \norm{u_0-1}+\norm{v-1}\leq Ct
\]
such that
\[
 M=M_0\barotimes D,\qquad
 N=v\bigl((u_0N_0u_0^*)\barotimes D\bigr)v^*
\]
form a prepared pair and $\dcb(M,N)\leq Ct$.
\end{proposition}

\begin{proof}
A countably decomposable infinite factor is properly infinite and contains
a unital copy of $\B(\ell^2)$: split its identity into countably many
equivalent projections and choose corresponding matrix units
\cite[III.1.3.5 and III.1.5.1, pp.~243, 247]{BlackadarOperatorAlgebras}.
\Cref{thm:amenable} embeds this amenable subfactor of $M_0$ into $N_0$
by a unitary of norm distance at most $150t$ from 1.  Conjugating $N_0$
in the inverse direction gives $u_0$ for which the subfactor is common.
This also proves that $N_0$ is infinite.
The adjusted ordinary gap is at most $301t$.
\Cref{lem:common-infinite-amplification} makes this also a cb gap.
\Cref{lem:tensor-gap} transfers it to the tensor pair because $D'$ is
hyperfinite.  By \Cref{lem:fixed-tensor-factor}, both tensor factors are
of type $\III_1$ and absorb $R_\lambda$.

For the moment write $N_1=(u_0N_0u_0^*)\barotimes D$.  The
bicentralizer and expected-subfactor results just stated give an
irreducible hyperfinite $\II_1$ factor $R\subseteq M$ and a faithful
normal state $\varphi$ on $M$ centralizing $R$.
The given representation of $M$ has a cyclic separating vector realizing
$\varphi$.  Here is the representation point needed for this assertion.
A normal representation on a separable space splits as a countable sum
of cyclic normal representations.  Each cyclic summand embeds in the
standard representation by the standard-form vector implementing its
vector functional \cite[Lemma~2.10(1), p.~278]{Haagerup1975}.
A countable amplification of the standard
representation is unitarily equivalent to it, since its type $\III$
commutant contains a countable family of isometries with ranges summing
to 1.  Finally, every nonzero projection in that countably decomposable
type $\III$ commutant is equivalent to 1
\cite[Theorem~III.1.7.9]{BlackadarOperatorAlgebras}.  Thus every nonzero normal
representation here is equivalent to the standard one.  Transport the
standard vector for $\varphi$ to obtain the desired $\Omega$ on the
actual space $H\otimes K$.  No conjugation of either algebra is made in
this choice.

Put $S=J_MRJ_M\subseteq M'$.  Centralization gives
\begin{equation}
 r\Omega=J_Mr^*J_M\Omega\qquad(r\in R).
 \label{eq:common-central-vector}
\end{equation}
First apply \Cref{thm:amenable} to $R\subseteq M$ and its near
inclusion in $N_1$.  It gives a unitary $v_1$ with
$v_1Rv_1^*\subseteq N_1$ and $\norm{v_1-1}=O(t)$.  Put
$N_2=v_1^*N_1v_1$, so $R\subseteq M\cap N_2$.
Cb gaps change by at most twice the norm distance of the adjusting
unitary from 1.  Thus \Cref{lem:commutant-cb} gives
$d(M',N_2')=O(t)$.  Apply \Cref{thm:amenable} to $S\subseteq M'$
and its near inclusion in $N_2'$.  It gives a unitary $v_2$ with
\[
 v_2Sv_2^*\subseteq N_2',\qquad
 \norm{v_2-1}=O(t),\qquad
 v_2\in(S\cup N_2')''\subseteq R'.
\]
Set $N=v_2^*N_2v_2$.  Since $v_2$ commutes with $R$, this algebra
still contains $R$, and now $S\subseteq N'$.
The algebra $M$, its state $\varphi$, and the vector $\Omega$ have
remained fixed.  The unitary in the statement is $v=v_2^*v_1^*$,
with $\norm{v-1}=O(t)$, and the resulting cb gap is $O(t)$.

We next check irreducibility and the vector, rather than assume they
survive an arbitrary small perturbation.  Choose increasing matrix
subfactors $R_j$ strongly generating $R$.  Haar averaging over
$\U(R_j)$ converges weakly on $M$ to $x\mapsto\varphi(x)1$.
Indeed every bounded cluster point commutes with all $R_j$, hence is
scalar, and the averages preserve $\varphi$.  If a nontrivial projection
$p\in R'\cap N$ existed, a close approximant $x\in M$ would therefore
give
\[
 \dist(p,\C1)\leq\norm{p-x}=O(t).
\]
The left side is $1/2$.  Taking $t$ sufficiently small proves
$R'\cap N=\C1$.  The same argument in the commutants, with $S$ and
the vector state on $M'$, proves $S'\cap N'=\C1$.

Let $p$ be the projection onto $\overline{N'\Omega}$.  It belongs to
$N$.  By \eqref{eq:common-central-vector}, this subspace reduces $R$:
for $r\in R$ and $x'\in N'$, the vector $rx'\Omega=x'r\Omega$ lies
in $\overline{N'\Omega}$, and the same holds for $r^*$.
Thus $p\in R'\cap N$, so $p=1$.  Similarly the projection onto
$\overline{N\Omega}$ belongs to $S'\cap N'$ and is 1.  Hence
$\Omega$ is cyclic and separating for $N$.  The common commuting
right action in \eqref{eq:common-central-vector} shows that its state
on $N$ is invariant under conjugation by $\U(R)$, and hence centralizes
$R$.  The two operators $J_MrJ_M,J_NrJ_N\in N'$ now agree on
$\Omega$; since $\Omega$ is separating for $N'$, they agree everywhere.
All conditions in \Cref{def:prepared-pair} are satisfied.
\end{proof}

\subsection{Flattening rescaled Tomita graphs}

Fix a prepared pair on a space denoted again by $H$. Choose increasing
unital matrix subfactors $R_j$ strongly generating its common subfactor $R$.
For $P=M,N$,
let $S_P$ be the closure of $x\Omega\mapsto x^*\Omega$ and write
$S_P=J_P\Delta_P^{1/2}$.  For $c>0$, let $E_c^P$ be the orthogonal
projection in $H\oplus\overline H$ onto the closure of
\begin{equation}
 g_c(x)=(x\Omega,c\,\overline{x^*\Omega}),\qquad x\in P.
 \label{eq:rescaled-graph}
\end{equation}
The conjugate Hilbert space makes this a complex linear graph.
Our next result gives uniform control even as $c$ tends to zero or
infinity.  Bounding $x$ in the ordinary operator norm directly would
lose this uniformity.  We instead flatten a graph vector and then
replicate it into a row, where cb closeness supplies the needed estimate.

\begin{proposition}[Uniform graph comparison]
\label{prop:modular-graphs}
There is an absolute constant $C$ such that every prepared pair with
$\dcb(M,N)\leq\gamma$ satisfies
\[
 \sup_{c>0}\norm{E_c^M-E_c^N}\leq C\gamma.
\]
\end{proposition}

\begin{proof}
For $s\in R$ put $R_s=J_Ms^*J_M=J_Ns^*J_N$.  On either graph,
the operation $g_c(x)\mapsto g_c(rxs)$ for $r,s\in\U(R)$ is the
restriction of the common unitary
\[
 (rR_s)\oplus\overline{s^*R_{r^*}}
 \quad\text{on }H\oplus\overline H.
\]
Both graph subspaces reduce these unitaries.  Thus
$T=(1-E_c^N)E_c^M$ commutes with them.

\paragraph{A bounded element retaining the graph signal.}
First take $c=\sqrt{k}$, $k\in\N$.  If $T\ne0$, choose $x\in M$
with
\[
 \norm{g_c(x)}=1,
 \qquad\norm{Tg_c(x)}\geq\tfrac12\norm{T}.
\]
Such vectors are dense in the graph.  Write
\[
 \alpha=\varphi(x^*x),\qquad\beta=\varphi(xx^*),
 \qquad\alpha+k\beta=1.
\]
Choose independent Haar unitaries $u_i,v_i\in\U(R_j)$ and independent
signs $\epsilon_i\in\{-1,1\}$, and put
\[
 x_i=u_ixv_i,\qquad y=m^{-1/2}\sum_{i=1}^m\epsilon_ix_i.
\]
For fixed unitaries, the vectors $Tg_c(x_i)$ all have the same norm.
The sign second moment and the elementary Hilbert fourth-moment bound
give
\[
 \E_\epsilon\norm{Tg_c(y)}^2=\norm{Tg_c(x)}^2,
 \qquad
 \E_\epsilon\norm{Tg_c(y)}^4\leq3\norm{Tg_c(x)}^4.
\]
Consequently, with sign probability at least $1/12$,
\begin{equation}
 \norm{Tg_c(y)}\geq\frac{\norm{T}}{2\sqrt2}.
 \label{eq:graph-sign-signal}
\end{equation}
For completeness, this probability bound follows by applying
Cauchy--Schwarz to the part of a nonnegative random variable above half
its mean: $\mathbb P(Z\geq\E Z/2)\geq(\E Z)^2/(4\E Z^2)$.

We need a simultaneous bound on the fourth moments of $y$.  The Haar
averages
\[
 \mathcal E_j(a)=\int_{\U(R_j)}uau^*\,du\qquad(a\in M)
\]
converge weakly to $\varphi(a)1$ as in the preceding proof.  For
distinct $i,l$, centralizer cycling and Haar invariance give the exact
identities
\begin{align*}
 \E\varphi(x_i^*x_i x_l^*x_l)
   &=\varphi\bigl((x^*x)\mathcal E_j(x^*x)\bigr),\\
 \E\varphi(x_i^*x_l x_l^*x_i)
   &=\varphi\bigl(x^*\mathcal E_j(xx^*)x\bigr).
\end{align*}
For the first identity, cycle $v_i^*$ to the right and average over
the Haar unitary $v_iv_l^*$.  For the second, cycle $v_i^*$ and $v_i$
away and average over $u_i^*u_l$.  Each formula contains just one
weakly convergent average.  The three sign-pairing terms in
$\varphi((y^*y)^2)$ have expected limits
\begin{align*}
 \E\varphi(x_i^*x_i x_l^*x_l)&\longrightarrow\alpha^2,\\
 \E\varphi(x_i^*x_l x_l^*x_i)&\longrightarrow\alpha\beta,\\
 \abs{\E\varphi(x_i^*x_l x_i^*x_l)}
 &\leq\E\varphi(x_i^*x_l x_l^*x_i)
 \longrightarrow\alpha\beta.
\end{align*}
For the third, state Cauchy--Schwarz gives
$\abs{\varphi(a^2)}\leq\varphi(aa^*)^{1/2}\varphi(a^*a)^{1/2}$,
and the two expected factors agree by interchanging $i,l$.
The all-equal indices contribute $m^{-1}\varphi((x^*x)^2)$.
The same calculation for $yy^*$ interchanges $\alpha,\beta$.
For the fixed $x,k$, first choose $m$ so that the weighted diagonal
contribution satisfies
\[
 \frac{\varphi((x^*x)^2)+k\varphi((xx^*)^2)}{m}\leq1.
\]
The limiting weighted off-diagonal contribution is at most
$\alpha^2+k\beta^2+2(k+1)\alpha\beta\leq6$, because
$\alpha+k\beta=1$ and $k\geq1$.  Now choose $j$ large enough that
the weighted off-diagonal contribution is at most 7.
Since $\E\norm{g_c(y)}^2=1$, these choices give
\begin{equation}
 \E\left[\norm{g_c(y)}^2+
       \varphi((y^*y)^2)+k\varphi((yy^*)^2)\right]\leq10.
 \label{eq:graph-first-moment-bound}
\end{equation}
Markov's inequality bounds the probability that the bracket in
\eqref{eq:graph-first-moment-bound} exceeds $240$ by $1/24$.
Together with \eqref{eq:graph-sign-signal}, this gives a realization
with both properties.  Fix such a $y$, and clip its singular values at
a fixed constant $K_1$:
\[
 z=y f_{K_1}(|y|),\qquad
 f_b(t)=\min(1,b/t),\quad f_b(0)=1.
\]
Functional calculus on both sides of the polar decomposition gives
\[
 \norm{g_c(y-z)}^2
 \leq K_1^{-2}
 \left[\varphi((y^*y)^2)+k\varphi((yy^*)^2)\right].
\]
Choose $K_1$ once so that $\sqrt{240}/K_1\leq1/(4\sqrt2)$.
Writing $a_0=1/(4\sqrt2)$, the signal loss satisfies
\[
 \begin{aligned}
 \norm{Tg_c(z)}
 &\geq\norm{Tg_c(y)}-\norm{T}\norm{g_c(y-z)}\\
 &\geq\left(\frac1{2\sqrt2}-\frac{\sqrt{240}}{K_1}\right)\norm{T}
 \geq a_0\norm{T}.
 \end{aligned}
\]
Thus the choice of $K_1$ is independent of the size of $\norm{T}$, and
\begin{equation}
 \norm{z}\leq K_1,\quad
 \norm{Tg_c(z)}\geq a_0\norm{T},\quad
 \varphi(z^*z)\leq240,\quad\varphi(zz^*)\leq240/k.
 \label{eq:graph-clipped-vector}
\end{equation}

\paragraph{A bounded row retaining the normalized signal.}
We have produced a bounded element retaining the signal, but a single
bounded-element approximation would still cost $\sqrt{k}\gamma$ on
the second graph leg.  To remove that factor, replicate $z$ into the row
\[
 Y=[r_1z\ \cdots\ r_kz],\qquad A=YY^*,
\]
where the $r_i$ are independent Haar unitaries in a sufficiently large
$R_j$.  For rows of graph vectors use the normalized Hilbert norm
\[
 \norm{(\eta_i)_{i=1}^k}_{\mathrm{av}}^2
       =k^{-1}\sum_i\norm{\eta_i}^2.
\]
Equivariance makes the normalized signal of this row at least
$a_0\norm{T}$ for every choice of the $r_i$.
Put $\alpha_z=\varphi(z^*z)$ and $\beta_z=\varphi(zz^*)$.
Haar averaging and centralizer cycling give
\begin{align*}
 \E\varphi(A^2)&\longrightarrow
 k\varphi((zz^*)^2)+k(k-1)\beta_z^2,\\
 \E\left[k^{-1}\sum_i\varphi(Y_i^*AY_i)\right]
 &\longrightarrow
 \varphi((z^*z)^2)+(k-1)\alpha_z\beta_z.
\end{align*}
For example, the second off-diagonal term averages $zz^*$ under
$r_i^*r_l$ between $z^*$ and $z$.  By
\eqref{eq:graph-clipped-vector}, the limits are bounded by an absolute
constant: use $\varphi((zz^*)^2)\leq K_1^2\beta_z$ and
$\varphi((z^*z)^2)\leq K_1^2\alpha_z$ for the diagonals.
Choose the stage and then a realization with
\begin{equation}
 \varphi(A^2)+k^{-1}\sum_i\varphi(Y_i^*AY_i)\leq C_1,
 \label{eq:graph-row-moments}
\end{equation}
where $C_1$ is absolute.

Clip the whole row on the left:
\[
 X=f_{K_2}(A^{1/2})Y,
\]
so $\norm{X}\leq K_2$.  If $Q=1-f_{K_2}(A^{1/2})$, then
$Q^2\leq A/K_2^2$ and $Y-X=QY$.  As $c^2=k$, its normalized squared
graph error is
\begin{align*}
 k^{-1}\sum_i\norm{g_c(Y_i-X_i)}^2
 &=k^{-1}\sum_i\varphi(Y_i^*Q^2Y_i)+\varphi(QAQ)\\
 &\leq C_1/K_2^2.
\end{align*}
Choose the absolute $K_2$ so that $\sqrt{C_1}/K_2\leq a_0/2$.
The same proportional error estimate for the row gives
\begin{equation}
 \begin{aligned}
 \norm{(Tg_c(X_i))_i}_{\mathrm{av}}
 &\geq\norm{(Tg_c(Y_i))_i}_{\mathrm{av}}
       -\norm{T}\norm{(g_c(Y_i-X_i))_i}_{\mathrm{av}}\\
 &\geq\left(a_0-\frac{\sqrt{C_1}}{K_2}\right)\norm{T}
 \geq(a_0/2)\norm{T}.
 \end{aligned}
 \label{eq:graph-final-signal}
\end{equation}

\paragraph{Comparison at one matrix level and all rescalings.}
Cb closeness gives a row $Z\in M_{1,k}(N)$ with
$\norm{X-Z}\leq K_2\gamma$, with arbitrarily small slack if needed.
For any row $[e_i]$ of norm at most $\eta$,
\[
 k^{-1}\sum_i\norm{g_c(e_i)}^2
 =k^{-1}\sum_i\varphi(e_i^*e_i)+
   \varphi\left(\sum_i e_ie_i^*\right)
 \leq2\eta^2.
\]
Here the graph-vector formula itself makes sense for all operators on
$H$, whether or not they lie in $M$ or $N$.
Because $g_c(X_i)\in\operatorname{ran}E_c^M$ and
$g_c(Z_i)\in\operatorname{ran}E_c^N$,
\[
 \norm{(Tg_c(X_i))_i}_{\mathrm{av}}
 =\norm{((1-E_c^N)g_c(X_i))_i}_{\mathrm{av}}
 \leq\sqrt2K_2\gamma.
\]
Together with \eqref{eq:graph-final-signal}, this bounds
$\norm{(1-E_c^N)E_c^M}$ by an absolute multiple of $\gamma$.
If $T=0$ the same assertion is immediate.  Interchanging $M,N$ gives
the reverse bound, hence the claimed projection bound for $c=\sqrt{k}$.

The antiunitary flip
$(\xi,\overline\eta)\mapsto(\eta,\overline\xi)$ takes
$g_c(x)$ to $c g_{1/c}(x^*)$, so the bound also holds for
$c=1/\sqrt{k}$.  For $c\geq1$ choose $k=\lfloor c^2\rfloor$;
then $c/\sqrt{k}\in[1,\sqrt2]$.  Scaling the second leg by this ratio
has condition number at most $\sqrt2$ and increases either one-sided
subspace gap by at most that factor.  Reciprocal scaling handles
$0<c<1$.  Finally, for two orthogonal projections $E,F$,
\[
 \norm{E-F}=\max\{\norm{(1-F)E},\norm{(1-E)F}\}.
\]
This proves the uniform assertion.
\end{proof}

\subsection{Separated spectral cuts}

\begin{corollary}[Logarithmic modular cuts]
\label{cor:modular-cuts}
For a prepared pair with $\dcb(M,N)\leq\gamma$, put
$L_P=\log\Delta_P$, $P=M,N$.  There is an absolute constant $C$ such
that, for every $a\in\R$,
\[
 \norm{1_{(-\infty,a-1/2]}(L_N)
        1_{[a+1/2,\infty)}(L_M)}\leq C\gamma.
\]
The same estimate holds with $M,N$ interchanged.
\end{corollary}

\begin{proof}
The upper-left block of the projection onto a closed operator graph is
$(1+A^*A)^{-1}$.  Applied to \eqref{eq:rescaled-graph}, this gives
$(1+c^2\Delta_P)^{-1}$.  With $g(s)=(1+e^s)^{-1}$,
\Cref{prop:modular-graphs} consequently yields
\[
 \norm{g(L_M-a)-g(L_N-a)}\leq C\gamma\qquad(a\in\R).
\]
Let $p=1_{(-\infty,a-1/2]}(L_N)$ and
$q=1_{[a+1/2,\infty)}(L_M)$, and regard $pq$ as an operator
$qH\to pH$.  Let $A$ and $B$ be the restrictions of $g(L_N-a)$ and
$g(L_M-a)$ to those respective subspaces.  Then
\[
 A\geq g(-1/2)1,\qquad B\leq g(1/2)1,
 \qquad\norm{A(pq)-(pq)B}\leq C\gamma.
\]
For $D=A(pq)-(pq)B$, integration of the derivative of
$e^{-tA}(pq)e^{tB}$ gives
\[
 pq=\int_0^\infty e^{-tA}D e^{tB}\,dt,
 \qquad
 \norm{pq}\leq\frac{\norm{D}}{g(-1/2)-g(1/2)}.
\]
The denominator is a fixed positive number.  Reversing the two algebras
proves the other estimate.
\end{proof}
